\section{Proposed solution}
\label{sec:solution}

\subsection{Problem statement}

The project asks for a personal assistant that hears the user, reads the
user's face, and answers questions about a set of machine learning problems.
Its scope has four deliverables: an agent that recognises the emotion of the
person from a live camera frame, an agent that turns speech into text and
executes the instruction, at least ten machine learning models chosen from
the list of the course, and a set of spoken commands that run those models.
The models are served by a REST API, and the user talks to them through a
desktop application.

\jarvis{} is the answer to that statement. It is a Spanish-speaking desktop
assistant in which every request follows one path: the microphone records
the user, a speech-to-text model transcribes the audio, an intent parser
decides which of the ten models the sentence refers to and which values it
carries, the REST API runs that model, and the reply is written and spoken
back in a tone that depends on the emotion the webcam read on the user's
face. An animated penguin avatar changes its expression with that emotion.

\subsection{Design decisions}

Three decisions shape the solution and are defended here rather than in the
code.

\paragraph{Everything runs locally.}
The statement allows cloud services from Google or Azure for face and speech
recognition in the first delivery. \jarvis{} does not depend on them: the
emotion classifier is a convolutional network exported to ONNX and run with
\code{onnxruntime} on the CPU, speech-to-text is Whisper
through \code{faster-whisper}, and text-to-speech is Piper with a Mexican
Spanish voice. The only optional cloud call is to Gemini for intent parsing,
and when no key is configured a rule-based parser takes over with the same
contract, so the assistant works offline end to end. This makes the demo
independent of quotas, network and billing, and keeps the user's camera
frames and voice on the machine.

\paragraph{Ten models, one contract.}
Each model is a CRISP-DM notebook that ends by writing a
\code{joblib} bundle with the estimator and everything the API needs to feed
it without retraining: the predictor list, encodings, fill values for
missing inputs, thresholds, class names and training ranges. The API loads
the bundles once at startup and exposes them behind one endpoint shape,
\code{POST /predict/\{model\_name\}}. The application never touches an
estimator; it only speaks JSON. Adding an eleventh model is a notebook, an
adapter and a schema, and the assistant discovers the new schema at startup.

\paragraph{The emotion changes how the answer is said, not what it says.}
The classifier returns one of seven emotions (angry, disgust, fear, happy,
sad, surprise, neutral), the vocabulary of the FER-2013 challenge, which in turn follows Ekman's basic emotions. The prediction of the model is computed first and is
identical whatever the face shows; a tone table then wraps it in a prefix
and a suffix (calm and short for an angry user, cheerful for a happy one).
This keeps the machine learning answer auditable and confines the affective
layer to phrasing.

\subsection{Commands}

The assistant wakes on the word \emph{JarvisTEC} (and the mishearings
Whisper produces for it, such as \emph{Yarvis} or \emph{Yalvis}, all matched on
a word boundary). One sentence per model is enough when it carries every
value the model needs; when a required value is missing the assistant asks
for it and keeps the pending intent across turns. Three example sentences
per model, in Spanish, ship with the application in \code{app/prompts.txt};
Table~\ref{tab:commands} gives one per model.

\begin{table}[htbp]
  \centering
  \small
  \begin{tabularx}{\textwidth}{@{}lL@{}}
    \toprule
    Model & Example command (Spanish, as spoken) \\
    \midrule
    Bitcoin & \emph{JarvisTEC, ¿cuánto va a costar el bitcoin mañana?} \\
    Car price & \emph{JarvisTEC, ¿cuánto vale un carro diesel del 2015 que nuevo costó 10 lakhs, tiene 40 mil kilómetros, un solo dueño, lo vende una agencia y es manual?} \\
    Churn & \emph{JarvisTEC, ¿se va a ir un cliente hombre \ldots contrato mes a mes, factura electrónica, paga con cheque electrónico, lleva 3 meses y paga 95 dólares al mes?} \\
    Wine & \emph{JarvisTEC, ¿qué calidad tiene un vino blanco con 13.5 de alcohol, densidad 0.99, \ldots y 2 de azúcar residual?} \\
    Flight delay & \emph{JarvisTEC, ¿cuánto se va a retrasar el vuelo de American Airlines de Dallas a Chicago el lunes 15 de enero a las 12 del mediodía, con 802 millas?} \\
    Stroke & \emph{JarvisTEC, ¿qué riesgo de derrame tiene un hombre de 68 años, hipertenso, \ldots con glucosa de 140, índice de masa corporal 31 y que dejó de fumar?} \\
    Cirrhosis & \emph{JarvisTEC, ¿tiene cirrosis una mujer de 55 años sin edema, sin ascitis, con hepatomegalia, \ldots plaquetas 200 y protrombina 11?} \\
    Avocado & \emph{JarvisTEC, ¿a cuánto estaría el aguacate orgánico en San Francisco el 10 de septiembre de 2017 con un volumen total de 20 mil \ldots?} \\
    Hepatitis C & \emph{JarvisTEC, ¿qué categoría de hepatitis tiene un hombre de 45 años con albúmina 38, fosfatasa alcalina 60, ALT 30, AST 90, \ldots y proteínas 70?} \\
    Body fat & \emph{JarvisTEC, ¿cuál es mi porcentaje de grasa corporal si mi abdomen mide 90 centímetros, peso 175 libras y mi muñeca mide 18 centímetros?} \\
    \bottomrule
  \end{tabularx}
  \caption{One spoken command per model. The full set, three per model, is in \code{app/prompts.txt}.}
  \label{tab:commands}
\end{table}

\subsection{User interface}

The window (Figure~\ref{fig:ui}) has the avatar and the camera card on the
left and the conversation on the right. The camera card shows the mirrored
webcam feed with the detected face boxed and the emotion and its confidence
written under it. The input row has a single button that is the microphone
while the entry is empty, the send arrow once there is text, and a red
recording indicator while the microphone is open. A status strip at the
bottom reports the microphone, the emotion detector, the API, the
speech-to-text engine and which parser answered (a Gemini model or the
rules). System notices never enter the chat.

\begin{figure}[htbp]
  \centering
  \begin{tikzpicture}[
      every node/.style={font=\small},
      card/.style={draw, rounded corners, text width=3.6cm, align=center},
    ]
    \node[card, minimum height=3.2cm] (avatar) {Penguin avatar\\ \scriptsize idle / listening / thinking / speaking, one expression per emotion};
    \node[card, minimum height=1.6cm, below=0.3cm of avatar] (camera) {Camera card\\ \scriptsize face box + emotion label, ON/OFF};
    \node[card, text width=4.6cm, minimum height=3.9cm, anchor=north west] (chat) at ([xshift=0.5cm]avatar.north east) {Conversation\\ \scriptsize one bubble per turn};
    \node[card, text width=4.6cm, minimum height=0.9cm, below=0.3cm of chat] (input) {Entry + MIC / send button};
    \node[card, text width=8.9cm, minimum height=0.8cm, anchor=north west] (status) at ([yshift=-0.3cm]camera.south west) {Status strip: MIC \; EMO \; API \; STT \; IA};
  \end{tikzpicture}
  \caption{Layout of the \jarvis{} window. \code{Enter} sends typed text, \code{Ctrl+Space} toggles recording, \code{Ctrl+E} toggles the webcam, \code{Esc} quits.}
  \label{fig:ui}
\end{figure}
